% sample_with_error.tex — 包含多种典型错误的 LaTeX 测试样本
% 用于测试 LatexErrorParser 的错误提取与分类能力

\documentclass[12pt,a4paper]{article}

% === 故意引入冲突的宏包 ===
\usepackage{fontspec}
\usepackage[utf8]{inputenc}   % fontspec 下不需要 inputenc
\usepackage{fontenc}          % 与 fontspec 冲突

% === 故意使用不存在的字体（触发 FONT_NOT_FOUND）===
\setmainfont{完全不存在的字体名称ABCXYZ}

% === 重复加载宏包 ===
\usepackage{amsmath}
\usepackage{amsmath}

% === 缺失宏包（正文引用但未加载）===
% 注意: 未加载 hyperref

\newcommand{\testcmd}[1]{#1}

\begin{document}

\section{测试}
这是一个故意包含错误的文档。

% === 未定义命令 ===
\undefinedcommand{hello}

% === 数学模式错误 ===
\begin{equation}
    x = y +
\end{equation}

% === 引用不存在的标签 ===
在第~\ref{sec:nonexistent} 节中...

\section{另一个测试}

\end{document}
